\section{Discussion}

\subsection{What the Paradox Means for Sugar Tax Evidence}

Sugar-sweetened beverage taxes have been implemented or proposed in over 40 countries \cite{who2023ssb}, drawing in part on cross-sectional evidence linking sugar availability to metabolic disease \cite{basu2013, basu2013ajph}. Individual-level prospective cohorts and meta-analyses consistently link sugar-sweetened beverage consumption to type 2 diabetes and metabolic syndrome \cite{imamura2015, malik2010}, and a recent global analysis attributed 2.2 million new diabetes cases to these beverages \cite{laracastor2025}. That evidence remains compatible with our result. The failure lies in national food-supply identification: country-level supply data mix sugar exposure with correlated development trends. Policy arguments built on cross-country food supply associations therefore rest on weaker identification than commonly assumed.

South Africa's sugar-sweetened beverage tax, implemented in April 2018, illustrates the difficulty. South Africa's NCD-RisC obesity trajectory shows no deflection after 2018, continuing its pre-tax linear trend. But the NCD-RisC smoothing model imposes temporal continuity by design, so the absence of a visible break cannot be taken as evidence that the tax had no effect. The data most commonly used to study national-level food-obesity relationships cannot resolve the kind of sharp policy discontinuity that would provide clean within-country identification. Evaluating SSB taxes requires survey data collected before and after implementation \cite{colchero2016}, not modeled trajectories that smooth over exactly the variation researchers need.

\subsection{What the Paradox Means for Dietary Risk Assessment}

The Global Burden of Disease study estimates diet-attributable deaths using intake data that faces the same structural problem \cite{gbd2019}. GBD dietary risk estimates combine survey data with food supply proxies and model-based extrapolation. To the extent that these estimates rely on cross-country variation in food supply to calibrate risk, they inherit the confounding structure we document. Risk attributed to specific food categories based on national supply data should be interpreted with caution.

Food Balance Sheets are a supply instrument, not an intake instrument. They measure national availability: production plus imports minus exports minus non-food uses, divided by population. They omit distribution within countries, food waste, processing, and actual dietary intake. Del Gobbo et al.\ \cite{delgobbo2015} compared FAO supply estimates to individual-level dietary surveys across 113 countries and found that FAO data overestimates intake by 50--270\% for most food groups. Vonderschmidt et al.\ \cite{vonderschmidt2024} showed that FAOSTAT's 2014 methodology revision creates discontinuities that can reverse apparent trends in food supply series, a problem for any study combining pre- and post-revision data. The gap between supply-level availability and individual consumption may be large enough to obscure genuine diet-disease relationships, even when those relationships exist at the individual level.

The identification problem extends beyond individual food groups. Recent FAOSTAT-based studies use total energy and nutritive-value supplies \cite{siminiuc2025,delgadopertinez2026}; our comparator shows that aggregate energy and macronutrient totals follow the same trend structure in SSA. Cross-national ecological studies using other food classification systems face the same design problem: Monteiro et al.\ \cite{monteiro2018} found that ultra-processed food availability predicted obesity across 19 European countries, a design that mixes co-trending with causal signal. FAOSTAT remains useful for health research when the estimand is explicit. Cross-national food-supply comparisons map development gradients; stronger designs are needed for causal diet claims.

\subsection{Better Identification Strategies}

If cross-sectional and one-way fixed-effects designs are inadequate, three approaches might provide stronger evidence.

First, policy discontinuities. Sugar taxes, trade liberalization events, or agricultural policy changes create sharp temporal variation in specific food categories. Difference-in-differences or synthetic control designs exploiting these events could provide within-country identification that panel fixed effects cannot.

Second, subnational data. District- or province-level panels increase the number of units and the variation in food environments while holding national-level trends constant. Subnational data also allows analysis of food retail access, supermarket penetration, and ultra-processed food availability, measures closer to dietary exposure than national supply totals.

Third, individual-level longitudinal surveys. Panel studies tracking dietary intake and health outcomes eliminate the ecological fallacy inherent in country-level analysis. Such data exist for some high-income countries but remain largely absent in Sub-Saharan Africa.

\subsection{Limitations}

Four limitations bear on interpretation. First, the panel covers 37 of 48 SSA countries; 11 are absent from the comparable FAOSTAT release. Second, NCD-RisC smoothing compresses within-country variation, working against detecting within-country associations. The consistency of the null across six specifications, including long differences that compare only decade endpoints, argues this is not the sole explanation. Third, we analyze food supply, not food intake. The paradox could reflect a failure of FAOSTAT to capture dietary variation rather than a genuine absence of diet-obesity relationships. Fourth, our panel period (2010--2022) is short enough that slow-acting dietary effects could be missed by trend-removal methods.

\subsection{Reframing the Question}

The field should redirect the question toward structural conditions that drive divergent obesity trajectories across countries. Our evidence points to initial conditions and development pathways as the primary source of cross-country variation. Dietary-composition claims require designs that observe exposure and outcome changes within the same units.

This is not simply the ecological fallacy by another name. Robinson \cite{robinson1950} warned that group-level associations may not hold for individuals. The ``Australian Paradox'' \cite{barclay2011} showed that sugar intake and obesity can diverge within a single country over time. We show something more specific than either: a quantifiable mechanism that makes an entire class of study designs unreliable, regardless of the food variable, the region, or the analytical refinement applied. Trend domination absorbs 99.4\% of outcome variance. The soybean oil case study shows that standard robustness tests---clustered standard errors, leave-one-out exclusion, Oster bounds---cannot distinguish co-trending from causation. The problem is global, affecting 98\% of countries.

Cross-national food-supply studies have built a large evidence base on development gradients in diet and obesity. Our results locate the boundary of that evidence. These designs can show which countries sit together in the global food-system transition; causal dietary effects require within-unit shocks, subnational panels, or individual longitudinal data. Those data exist in fragments but not yet at the scale required. Until they do, strong cross-sectional food-obesity associations should be read as development structure unless they survive trend-removal diagnostics.
